\SetPageTheme{story}
\PageHeader{Lectia 6}{Predarea algoritmului de descompunere}{L6 • P1}

\begin{missionbox}[Algoritmul-cheie]
  Acum le legam pe toate: operatorii, conditiile si repetitia. Vrem un algoritm care sa citeasca un numar si sa il desfaca in cifre, una cate una.

  Acesta este momentul in care problema incepe sa semene cu felul in care calculatorul citeste de fapt informatia: nu dintr-o privire, ci pas cu pas, bucata cu bucata.
\end{missionbox}

\begin{conceptbox}[Schema completa]
\begin{lstlisting}[style=codequest]
cat timp (n != 0) executa
    uc = n % 10;
    // aici folosim cifra extrasa
    n = n / 10;
sfarsit cat timp
\end{lstlisting}
\end{conceptbox}

\begin{guidebox}[Primul rol concret]
  Daca in locul comentariului scriem \texttt{afiseaza uc}, atunci algoritmul va afisa cifrele numarului in ordine inversa.
\end{guidebox}

\newpage
\SetPageTheme{guided}
\PageHeader{Lectia 6}{Prima aplicatie: numararea cifrelor}{L6 • P2}

\begin{examplebox}[Algoritm complet]
\begin{lstlisting}[style=codequest]
int n, cnt = 0;
citeste n;

daca (n == 0) atunci
    cnt = 1;
altfel
    cat timp (n != 0) executa
        cnt = cnt + 1;
        n = n / 10;
    sfarsit cat timp
sfarsit daca

afiseaza cnt;
\end{lstlisting}
\end{examplebox}

\begin{guidebox}[Exercitii ghidate]
  1. De ce \texttt{cnt} porneste de la 0?

  2. De ce \texttt{0} este tratat separat?

  3. Cate iteratii are bucla pentru \texttt{n = 3408}?
  \AnswerSpace{5}
\end{guidebox}

\newpage
\SetPageTheme{guided}
\PageHeader{Lectia 6}{Semi-ghidat}{L6 • P3}

\begin{solobox}[Completeaza lipsurile]
\begin{lstlisting}[style=codequest]
int n, cnt = 0;
citeste n;

daca (n == 0) atunci
    cnt = ______;
altfel
    cat timp (__________) executa
        cnt = cnt + 1;
        n = ________;
    sfarsit cat timp
sfarsit daca
\end{lstlisting}
\end{solobox}

\begin{solobox}[Exercitiu independent scurt]
  Scrie in cuvinte ce face algoritmul de mai sus.
  \AnswerSpace{4}
\end{solobox}

\newpage
\SetPageTheme{challenge}
\PageHeader{Lectia 6}{Independent}{L6 • P4}

\begin{solobox}[Fisa independenta]
  1. Scrie singur algoritmul care afiseaza toate cifrele unui numar in ordine inversa.

  2. Arata pe scurt ce afiseaza pentru \texttt{n = 5021}.

  3. Explica de ce aceasta schema este fundamentala pentru tot ce urmeaza.
  \AnswerSpace{8}
\end{solobox}
